\chapter*{ВВЕДЕНИЕ}
\addcontentsline{toc}{chapter}{ВВЕДЕНИЕ}

Модели табличных данных могут терять качество при переносе между организациями, регионами, временными интервалами и условиями измерения. Высокое качество на случайно выделенной части исходной выборки само по себе не характеризует качество на новом домене.

Цель работы --- исследовать методы обучения ансамблей деревьев решений, устойчивых к сдвигу распределения, и экспериментально оценить их применимость при source-only выборе модели. Для достижения цели требуется:
\begin{enumerate}
\item определить виды сдвига и условия переноса модели;
\item систематизировать рассматриваемые методы и их реализации;
\item задать воспроизводимый протокол source/target разделения;
\item разделить сравнение настроенных моделей и сравнение отдельных механизмов;
\item проанализировать экспериментальные результаты, отрицательные наблюдения и ограничения.
\end{enumerate}

Работа носит экспериментально-аналитический характер. Разработка нового алгоритма не заявляется. Исследовательские утверждения связываются с конкретными реализациями, протоколами и экспериментальными артефактами.
